\documentclass[a4paper]{scrartcl}
\usepackage[utf8]{inputenc}
\usepackage[english]{babel}
\usepackage{graphicx}
% SVG support: converted to vector PDF by Inkscape at compile time (needs -shell-escape, see latexmkrc).
% inkscapelatex=false keeps the diagram's own text and fonts instead of re-typesetting them with LaTeX.
\usepackage[inkscapelatex=false]{svg}
\usepackage{lastpage}
\usepackage{pgf}
\usepackage{wrapfig}
\usepackage{fancyvrb}
\usepackage{listings}
\usepackage{xcolor}
% Style for SQL / PL/pgSQL code listings
\lstdefinestyle{sql}{
  language=SQL,
  morekeywords={RETURNS,TRIGGER,DECLARE,BEGIN,END,IF,THEN,RAISE,EXCEPTION,RETURN,LANGUAGE,FUNCTION,REPLACE,EXECUTE,BEFORE,EACH,ROW,INTEGER,LIMIT,INTO,INNER,JOIN},
  basicstyle=\ttfamily\footnotesize,
  keywordstyle=\color{blue!70!black}\bfseries,
  commentstyle=\color{green!45!black}\itshape,
  stringstyle=\color{red!65!black},
  frame=single,
  breaklines=true,
  columns=fullflexible,
  keepspaces=true,
  showstringspaces=false
}
\usepackage{fancyhdr}
\pagestyle{fancy}

% Create header and footer
\headheight 27pt
\pagestyle{fancyplain}
\lhead{\footnotesize{Data Storage Paradigms, IV1351}}
\chead{\footnotesize{Project task 1 report}}
\rhead{}
\lfoot{}
\cfoot{\thepage\ (\pageref{LastPage})}
\rfoot{}

\title{Data Storage Paradigms, IV1351}
\subtitle{Data Storage Paradigms, IV1351}
\author{}
\date{\today}

\begin{document}

\maketitle
\noindent\textbf{Project members:} \\ \hfill
[Mo Wang, mowang@ug.kth.se] \\ \hfill
% [Enter name, email] \\ \hfill
% [Enter name, email] \\ \hfill

\section*{Declaration:}

By submitting this assignment, it is hereby declared that all group members listed above have contributed to the solution. It is also declared that all project members fully understand all parts of the final solution and can explain it upon request.

It is furthermore declared that the solution below is a contribution by the project members only, and specifically that no part of the solution has been copied from any other source (except for lecture slides at the course IV1351), no part of the solution has been provided by someone not listed as a project member above, and no part of the solution has been generated by a system.

% \section*{Tips for Report Writing}
% \textbf{REMOVE THIS SECTION BEFORE SUBMITTING THE REPORT.}\\

% \noindent \textit{The target audience has exactly the same skills as the author, except they do not know anything at all about the specific application described in the report.} \\

% Consider the following:

% \begin{itemize}
%   \item \textbf{The report must be \textit{centered around the requirements}. Which are they (Introduction), how did you work to meet them (Method), what is the solution that meets them (Result), and how can you be sure they are met (Discussion). This is the IMRaD method.} The requirements on the Introduction, Method, Result and Discussion chapters are described below under each chapter.

%   \item Is spelling and grammar correct? Is spoken language avoided?

%   \item Does the report have a good structure with sections, subsections and paragraphs?

%   \item Is the text clarified with images and/or other figures, and with links to the code in your Git repository? Remember that all figures (images, tables, graphs, code listings, etc) shall be numbered and have a short explaining text.
% \end{itemize}

\section{Introduction}

% \textbf{This chapter tells \textit{what} are you going to do.} 
% Explain the task and the requirements on the solution. It's important to clearly state the requirements. \textit{Also specify which other student you worked with when solving the tasks, or if you worked alone.}

This report concerns the design of a course database system. The database system stores courses, course layout versions, teaching allocations, and employees at a university. A conceptual diagram, along with a logical-physical ER-diagram is created for modeling the database and specifying the database schema/layout, as well as queries for creating the database tables, inserting sample table data, and additional database trigger operations for complying with requirement constraints. Additionally, the project task is carried out alone.

The database is intended to store university courses and employees (teachers), which are described as below. The purpose is to generate reports and statistics from the given data, for instance calculating the total teaching cost of a course instance in a given time interval.

The course database requirements can be divided into different parts, each with their constraints:

\subsection{Course layout, instances and activities}
The first part concerns course layout and course instances. For each course, a unique course code is assigned, while the course name, ECTS credits (högskolepoäng or hp), the minimum and maximum number of students are also specified in the course layout. A course instance is a particular instance with its layout, given in a specified academic year. The academic year and study/academic period, labelled as P1, P2, P3, and P4, and the number of students in the instance, are specified in the course instance.

Each course instance contain a number of teaching activities, labelled as lecture, lab, tutorial, seminar, examination, administration and other. More teaching activity types could be added into the database, such as help sessions. The total number of required hours for each teaching activity is calculated differently.
For some teaching activities, such as lecture, lab, tutorial and seminar, the number of required hours can be calculated by multiplying the number planned hours with a multiplication factor of a given teaching activity. In other words, if the multiplication factor for lecture is 5, that means for every hour of holding lecture, 5 hours are required including preparation time.
For other teaching activities, the number of required hours are derived from the number of students or ECTS credits. The formulae for calculating the required hours of examination and admin are stated as below (taken from project description):

\begin{lstlisting}
Examination hour = 32 + 0.725 * #Students
Admin hours = 2 * HP + 28 + 0.2 * #Students
\end{lstlisting}

\subsection{Departments, people and employees}
The university has a number of employees (or teachers), with their contact details, job title and salary recorded in the database. The salary format can be given monthly. Employee salaries can be changed, thus it is necessary to version them. Employees are affiliated with different departments and are allocated to different teaching activities of course instances. Allocations can be proceeded depending on the given skill set of employees. The university contains multiple departments, each with a manager who is also an employee. Each employee can also have a supervisor, which is also another employee.

Employees can be allocated to different course instances, while a course instance can have multiple employees being allocated in. However, one teacher must not be allocated more than 4 different course instance in a same academic period.

\section{Literature Study}

% This chapter must prove that you collected sufficient knowledge before starting development, instead of just hacking away without knowing how to complete a task. State what you have read and briefly summarize what you have learned.

In terms of literature study, the relevant fundamental and theoretical knowledge for carrying out the project task. These includes the database relational model, SQL language, database normalization, conceptual models, logical and physical models. The relational model, SQL language and database normalization concepts are covered both in the course book and the lecture presentations, while the conceptual models, logical and physical models are mostly covered in Leif Lindbäck's videos which is more tied with the project task requirement. Also, the same concepts are covered in the coursebook, but are stated more generally.

\subsection{Relational database model}

The relational database model organizes data into relations (tables). Each relation (table) consists of rows (tuple or records) and columns (attributes). A relation has both an intentional state, which is the same as domain definition, and an extensional state, which describes the current data currently being stored in the table. A table can be considered as a relation, since each tuple exists in the table domain, which is a cross product of the domains of individual attributes.

The relational algebra and relational calculus provides different relation (table) operations, such as selecting tuples and projecting attributes of a relation. Additionally, it also states that relationships between tables can be established using primary and foreign keys, and can be considered by an equijoin between the two tables. A primary key uniquely identifies each record in a table, while a foreign key references a record in another table and maintains referential integrity.

The relational database model is an important key concept in terms of organizing the given data in terms of courses and employees since the database contains multiple interconnected entities, such as courses, course instances, employees, departments, and teaching activities.

\subsection{SQL}
Structured Query Language (SQL) is used to define, manipulate, query, and retrieve data in relational databases. Instead of specifying the query in relational calculus which looks mathematically verbose, SQL offers a query specification that resembles natural language. The language include multiple parts, such as Data Definition Language (DDL) which is used to create and alter database table schema, Data Query Language (DQL) which is used to retrieve data from tables using \texttt{SELECT} and Data Manipulation Language (DML) which is used to \texttt{INSERT}, \texttt{UPDATE} and \texttt{DELETE} records.

During creating database tables, database constraints on tables are used to further restrict the domain. The constraints includes for instance primary key constraints, foreign key constraints that specify actions during \texttt{ON UPDATE} and \texttt{ON DELETE}, \texttt{UNIQUE} constraints that allows only one unique value composition of one or more columns, \texttt{NOT NULL} constraints that rejects \texttt{NULL} value and \texttt{CHECK} constraints. These constraints helps database integrity and prevents invalid records from being stored.

Additionally, database triggers in PostgreSQL are also important for ensuring database integrity of rows that spans over different table records with more complex validation logic. A trigger is automatically executed on a database table operation such as \texttt{INSERT}, \texttt{UPDATE}, or \texttt{DELETE}. If the validation fails for the database trigger, the operation is rejected.

\subsection{Database normalization}

Database normalization is a process to reduce redundancy and avoid data anomalies. The process concerns dividing a large anormal table into smaller tables with relationship between them that complies with the normal form. The main normal forms includes First Normal Form (1NF), Second Normal Form (2NF), and Third Normal Form (3NF).

First Normal Form requires all attributes to contain atomic values, such as \texttt{INT}, \texttt{FLOAT}, \texttt{VARCHAR}, or \texttt{DATE}, and all rows and columns must be unique. Second Normal Form requires First Normal Form and that non-key attributes (attributes that are not a part of candidate key), must depend on the entire candidate key than only a part of it. Third Normal Form specifies that non-key attributes must depend only on the candidate key, along with Second Normal Form, which means that they do not depend transitively on candidate keys.

Database normalization is also another important concept because it helps to avoid unnecessary duplication in the database and makes data maintainance easier, especially in database storage with different types of related information.

\subsection{Conceptual Data Modeling}
---

Conceptual Data Modeling

Conceptual data modeling describes the information requirements of a system at a high level, independently of a specific database management system or implementation. Entity-Relationship (ER) modeling is commonly used for this purpose. It identifies entities, their attributes, and the relationships between them. Relationships can also have cardinalities, describing how many instances of one entity can be associated with instances of another entity.

The conceptual modeling material, particularly the instructional videos provided for the project, helped develop an understanding of how textual requirements can be translated into an entity-relationship model. This is an important initial step because it allows the relationships and constraints of the application to be identified before considering the actual database implementation.

---

\subsection{Logical and Physical Database Design}
---

Logical and Physical Database Design

Logical database design translates the conceptual model into a relational structure. Entities are represented as tables, attributes become columns, and relationships are implemented using primary and foreign keys. At this stage, the structure is further refined to satisfy relational database principles and normalization requirements.

Physical database design concerns how the logical structure is implemented in a particular database management system. This includes defining data types, constraints, indexes, and other database-specific implementation details.

Studying these two design levels provided an understanding of how to progress from an abstract representation of the application to an actual relational database schema. The distinction is important because the conceptual model focuses on the information requirements, whereas the logical and physical models describe how those requirements can be represented and implemented in a database.

Overall, the literature study provided the necessary theoretical foundation for analyzing the project requirements, designing the database models, and implementing the database using SQL and PostgreSQL.

---

\section{Method}

\textbf{This chapter tells \textit{how} you solved the task.}

Explain how you worked when solving the tasks and how you evaluated that your solution met the requirements. \textit{Do not explain your solution and do not refer to code}, that belongs to the \textit{Result} chapter. More specific instructions for the content can be found under each task on the Project page in Canvas.

\section{Result}

\textbf{This chapter explains \textit{the result} of what you did.}

\textbf{The report must show that you have done the work yourself and that you have understood what you have done}, both of these goals are met by carefully explaining your solution here in the result chapter, and proving that it meets the requirements. \textit{State each requirement that is met} and explain \textit{how you met it}. Also include links to your code in your Git repository, and include also diagrams, see Figure \ref{fig:diag}, and other figures to illustrate your reasoning. All figures must be referenced in the text. Ask yourself if the solution is clearly explained, and if the reader will understand the application. What would you yourself want to know if you read about the application, is that included in the report? More specific instructions for the content can be found under each task on the Project page in Canvas. 

\begin{figure}[h!]
  \begin{center}
    \includesvg[width=\textwidth]{diagrams/conceptual}
    \caption{conceptual diagram}
    \label{fig:diag}
  \end{center}
\end{figure}

\begin{figure}[h!]
  \begin{center}
    \includesvg[width=\textwidth]{diagrams/logical-physical}
    \caption{Logical-physical ER-diagram}
    \label{fig:diag}
  \end{center}
\end{figure}

\begin{lstlisting}[style=sql]
CREATE OR REPLACE FUNCTION validate_4_instances_each_study_period()
RETURNS TRIGGER AS $$
DECLARE
    new_instance_id INTEGER;
    new_academic_period_id INTEGER;
    current_value INTEGER;
    max_value INTEGER;
BEGIN
    -- extract instance_id and academic_period_id
    SELECT course_instance.id, course_instance.academic_period_id
    INTO new_instance_id, new_academic_period_id
    FROM planned_activity
    INNER JOIN course_instance ON course_instance.id = planned_activity.course_instance_id
    WHERE planned_activity.id = NEW.planned_activity_id;

    -- count the number of teaching_allocation from given employee and academic period
    SELECT COUNT(DISTINCT course_instance.id)
    INTO current_value
    FROM teaching_allocation
    INNER JOIN planned_activity ON planned_activity.id = teaching_allocation.planned_activity_id
    INNER JOIN course_instance ON course_instance.id = planned_activity.course_instance_id
    WHERE teaching_allocation.employee_id = NEW.employee_id
      AND course_instance.academic_period_id = new_academic_period_id
      AND course_instance.id <> new_instance_id;

    SELECT max_courses
    INTO max_value
    FROM allocation_rule
    ORDER BY id
    LIMIT 1;

    IF current_value >= max_value THEN
        RAISE EXCEPTION 'Problem: A teacher cannot be assigned to more than % course instances each academic period. Employee id: %, already allocated: %', max_value, NEW.employee_id, current_value;
    END IF;

    RETURN NEW;
END;
$$ LANGUAGE plpgsql;

CREATE OR REPLACE TRIGGER trigger_4_instances_each_study_period
BEFORE INSERT OR UPDATE
ON teaching_allocation
FOR EACH ROW
EXECUTE FUNCTION validate_4_instances_each_study_period();
\end{lstlisting}

\pagebreak

\section{Discussion}

\textbf{This chapter \textit{analysis} the result presented in the previous section.} 

Evaluate your solution according to the assessment criteria found in the assessment-criteria documents, which are found under the bullet \textit{In the Discussion chapter of your report...}, under each task on the Project page in Canvas. You do not have to cover all specified criteria.

\section{Comments About the Course}

\textbf{This section is optional, but please at least write approximately how much time you spent on the assignment}, including lectures, labs, tutorials and seminars. This is of great help for course evaluation.

Also, any other comment(s) related to this course offering or to coming offerings is much appreciated. 


\end{document}
